%% ============================================================
%%  PLANTILLA DE APUNTES DE MATEMÁTICA
%%  Compilar con:  pdflatex plantilla_apuntes.tex   (2 veces)
%%  Para el índice analítico: makeindex plantilla_apuntes
%%  y luego pdflatex otra vez. (Con --shell-escape se hace solo.)
%% ============================================================
\documentclass[11pt,a4paper,oneside]{report}

%% ---------- Idioma y codificación ----------
\usepackage[T1]{fontenc}
\usepackage[utf8]{inputenc}
\usepackage{lmodern}
\usepackage[spanish,es-noshorthands]{babel}
\usepackage{microtype}

%% ---------- Página y párrafos ----------
\usepackage[margin=2.5cm]{geometry}
\usepackage{parskip}            % separación entre párrafos sin sangría
\usepackage{setspace}
\setstretch{1.1}
\usepackage{fancyhdr}
\pagestyle{fancy}
\fancyhf{}
\fancyhead[L]{\small\nouppercase{\leftmark}}
\fancyhead[R]{\small\thepage}
\renewcommand{\headrulewidth}{0.4pt}

%% ---------- Matemática ----------
\usepackage{amsmath,amssymb,amsthm,mathtools}
\usepackage{bm}

%% ---------- Figuras, tablas, colores ----------
\usepackage{graphicx}
\graphicspath{{figuras/}}       % pon tus imágenes en la carpeta "figuras"
\usepackage{xcolor}
\usepackage{booktabs}           % tablas bonitas (\toprule, \midrule)
\usepackage{array}
\usepackage{tabularx}
\usepackage{longtable}
\usepackage{caption}
\usepackage{float}              % permite [H] para fijar posición
\usepackage{mwe}                % solo para imágenes de ejemplo; bórralo luego
\usepackage{tikz}

%% ---------- Código ----------
\usepackage{listings}
\lstset{
  basicstyle=\ttfamily\small,
  keywordstyle=\color{blue!70!black}\bfseries,
  commentstyle=\color{green!40!black}\itshape,
  stringstyle=\color{orange!80!black},
  numbers=left, numberstyle=\tiny\color{gray}, numbersep=8pt,
  backgroundcolor=\color{gray!8},
  frame=single, rulecolor=\color{gray!40},
  breaklines=true, showstringspaces=false, tabsize=4,
  literate={á}{{\'a}}1 {é}{{\'e}}1 {í}{{\'i}}1 {ó}{{\'o}}1 {ú}{{\'u}}1
           {Á}{{\'A}}1 {É}{{\'E}}1 {Í}{{\'I}}1 {Ó}{{\'O}}1 {Ú}{{\'U}}1
           {ñ}{{\~n}}1 {Ñ}{{\~N}}1
}

%% ---------- Índices ----------
\usepackage{imakeidx}
\makeindex[title=Índice analítico, intoc]

%% ---------- Cajas de colores ----------
\usepackage[most]{tcolorbox}

%% ---------- Hipervínculos y referencias (SIEMPRE al final) ----------
\usepackage[colorlinks=true,linkcolor=blue!60!black,citecolor=green!40!black,
            urlcolor=blue!60!black]{hyperref}
\usepackage[spanish,nameinlink]{cleveref}

%% ============================================================
%%  ENTORNOS: definición, teorema, demostración, notas...
%%  Todos comparten UN solo contador, dentro de cada capítulo.
%% ============================================================
\theoremstyle{definition}
\newtheorem{definicion}{Definición}[chapter]
\newtheorem{ejemplo}[definicion]{Ejemplo}
\newtheorem{ejercicio}[definicion]{Ejercicio}

\theoremstyle{plain}
\newtheorem{teorema}[definicion]{Teorema}
\newtheorem{proposicion}[definicion]{Proposición}
\newtheorem{lema}[definicion]{Lema}
\newtheorem{corolario}[definicion]{Corolario}

\theoremstyle{remark}
\newtheorem*{observacion}{Observación}
\newtheorem*{notademo}{Nota sobre la demostración}
\newtheorem*{idea}{Idea clave}

\renewcommand{\proofname}{Demostración}
\renewcommand{\qedsymbol}{$\blacksquare$}

% Nombres para cleveref
\crefname{definicion}{definición}{definiciones}
\Crefname{definicion}{Definición}{Definiciones}
\crefname{teorema}{teorema}{teoremas}
\Crefname{teorema}{Teorema}{Teoremas}
\crefname{proposicion}{proposición}{proposiciones}
\Crefname{proposicion}{Proposición}{Proposiciones}
\crefname{lema}{lema}{lemas}
\Crefname{lema}{Lema}{Lemas}
\crefname{corolario}{corolario}{corolarios}
\Crefname{corolario}{Corolario}{Corolarios}
\crefname{ejemplo}{ejemplo}{ejemplos}
\Crefname{ejemplo}{Ejemplo}{Ejemplos}
\crefname{ejercicio}{ejercicio}{ejercicios}
\Crefname{ejercicio}{Ejercicio}{Ejercicios}
\crefname{equation}{ecuación}{ecuaciones}
\Crefname{equation}{Ecuación}{Ecuaciones}
\crefname{figure}{figura}{figuras}
\Crefname{figure}{Figura}{Figuras}
\crefname{table}{tabla}{tablas}
\Crefname{table}{Tabla}{Tablas}
\crefname{chapter}{capítulo}{capítulos}
\Crefname{chapter}{Capítulo}{Capítulos}
\crefname{section}{sección}{secciones}
\Crefname{section}{Sección}{Secciones}

% Estilo visual: barra lateral de color + fondo suave
\tcbset{
  cajabase/.style={enhanced, breakable, sharp corners, boxrule=0pt,
    frame hidden, left=8pt, right=6pt, top=4pt, bottom=4pt,
    before skip=10pt, after skip=10pt}
}
\tcolorboxenvironment{definicion}{cajabase,
  colback=blue!6,   borderline west={3pt}{0pt}{blue!65!black}}
\tcolorboxenvironment{teorema}{cajabase,
  colback=red!6,    borderline west={3pt}{0pt}{red!70!black}}
\tcolorboxenvironment{proposicion}{cajabase,
  colback=orange!7, borderline west={3pt}{0pt}{orange!80!black}}
\tcolorboxenvironment{lema}{cajabase,
  colback=orange!7, borderline west={3pt}{0pt}{orange!80!black}}
\tcolorboxenvironment{corolario}{cajabase,
  colback=orange!7, borderline west={3pt}{0pt}{orange!80!black}}
\tcolorboxenvironment{ejemplo}{cajabase,
  colback=green!5,  borderline west={3pt}{0pt}{green!45!black}}
\tcolorboxenvironment{ejercicio}{cajabase,
  colback=green!5,  borderline west={3pt}{0pt}{green!45!black}}
\tcolorboxenvironment{notademo}{cajabase,
  colback=yellow!14, borderline west={3pt}{0pt}{yellow!60!black}}
\tcolorboxenvironment{observacion}{cajabase,
  colback=gray!10,  borderline west={3pt}{0pt}{gray!70}}
\tcolorboxenvironment{idea}{cajabase,
  colback=purple!6, borderline west={3pt}{0pt}{purple!70}}

%% ============================================================
%%  ATAJOS MATEMÁTICOS
%% ============================================================
\newcommand{\N}{\mathbb{N}}
\newcommand{\Z}{\mathbb{Z}}
\newcommand{\Q}{\mathbb{Q}}
\newcommand{\R}{\mathbb{R}}
\newcommand{\C}{\mathbb{C}}
\newcommand{\eps}{\varepsilon}
\newcommand{\dd}{\mathrm{d}}
\newcommand{\dpar}[2]{\frac{\partial #1}{\partial #2}}
\DeclarePairedDelimiter{\abs}{\lvert}{\rvert}
\DeclarePairedDelimiter{\norm}{\lVert}{\rVert}
\DeclarePairedDelimiter{\pesc}{\langle}{\rangle}
\DeclarePairedDelimiter{\conj}{\{}{\}}
\DeclareMathOperator{\sen}{sen}
\DeclareMathOperator{\tg}{tg}
\DeclareMathOperator{\Img}{Im}
\DeclareMathOperator{\Ker}{Ker}
\DeclareMathOperator{\rg}{rg}

% Atajo para "término indexado": lo escribe en el texto Y lo añade al índice
\newcommand{\termino}[1]{\textbf{#1}\index{#1}}

%% ============================================================
%%  PORTADA
%% ============================================================
\title{Apuntes de Estadística}
\author{Práxedis Jiménez Ruvalcaba}
\date{\today}

\begin{document}

\maketitle
\tableofcontents

%% ============================================================
\part{Curso de Estadística}
%% ============================================================

\chapter{Conceptos Introductorios}
Estos son apuntes realizados en la construcción de un curso de estadística que impartí. Se tomaron apuntes del
\begin{itemize}
  \item Casella, Berger; Statistical Inference
  \item Wackerly, Mendenhall, Schaffer; Estadística Matemática con Aplicaciones
  \item Freund, Miller; Mathematical Statistics with Applications
\end{itemize}

El curso iba enfocado al desarrollo de formalidad y familiraridad teórico-práctica intuitiva a una audiencia 
con cierta distancia al medio, así que se omitieron demostraciones y ejercicios de varias secciones. Puede ser que
ejercicios se añadan a posteriori.\\
Otra nota importante es que poner las gráficas e imágenes es cierto trabajo en latex, así que tal vez se vayan
añadiendo, pero un ejercicio recomendado es el de dibujarlas a mano, buscarlas o generarlas con algún programa
de python o R.

\section{Conceptos introductorios}

Como el curso se tomaba desde cierta carencia de familiaridad con la matemática, fue necesario aclarar algunos
conceptos básicos. No se abordaron desde la formalidad propia de un curso de conjuntos o probabilidad, porque 
la idea era generar una familiaridad con el lenguaje y notación para secciones posteriores del curso.

\begin{definicion}[Conjunto]
  Colección o agrupación bien definida de objetos (llamados elementos) que comparten característica o
  propiedades comunes. Estos cuentan con ciertas propiedades:\\
  \textbf{Contención}
  \[A\subset B \iff x \in A \Rightarrow x \in B\]
  Se dice que A está 'contenido' en B, o que A es un subconjunto de B.\\
  \textbf{Igualdad}
  \[ A = B \iff A \subset B \wedge B \subset A \]
\end{definicion}

\begin{definicion}[Operaciones de Conjuntos].\\
  \textbf{Unión}
  \[ A \cup B = \{ x | x \in A \vee x \in B\} \]
  \textbf{Intersección}
  \[ A \cap B = \{ x | x \in A \wedge x \in B\} \]
  \textbf{Complemento}
  \[ A^{c} = \{ x | x \not \in A\} \]
\end{definicion}

\begin{definicion}[Función]
  Una función es una regla que asigna a cada elemnto de un conjunto inicial (Dominio)
  un único elemento de un conjunto final (codominio).
\end{definicion}

\begin{definicion}[Evento]
  Un evento es una colección de posibles resultados de un experimento.\\
  Usualmente $\Omega$ denota el espacio muestral, este sería el conjunto de todos los posibles resultados. 
  Puede ser numerable o no numerable. 
\end{definicion}
La manera en que popularmente se entiende a los espacios muestrales numerables y no-numerables. Es que los numerables
son espacios muestrales de variables cualitativas, mientras que el no-numerable es para variables cuantitativa. 

\begin{definicion}[Eventos Disjuntos]
  Dos eventos $A$ y $B$ son disjuntos (mutuamente exluyentes) si 
  $A \cap B = \emptyset $. $A_{1},A_{2},...$ son disjuntos o 
  mutuamente exclusivos si $A_{i} \cap A_{j} = \emptyset$ $\forall i \not = j$
\end{definicion}

\begin{definicion}[Partición]
  Si $A_{1},A_{2},...$ son disjuntos y $\cup_{i=1}^{\infty}A_{i}=\Omega$, entonces
  $A_{1},A_{2},...$ forma una partición de $\Omega$
\end{definicion}

Las particiones son útiles para dividir el espacio muestral de forma de crear resultados interpretables o manejables.
Muchas de las categorías de datos continuos son ejemplos claros de partición, pudiera pensarse en calificaciones que
van de 0-100 donde los valores menores a 70 son reprobados y el resto son aprobados. De tener que hacer un análisis 
con valores potencialmente infinitos, se puede hacer con solo dos categorías. Se vuelve más interpretable o permite
trabajar de forma distinta, ahora desde un enfoque de las frecuencias en lugar de formalmente densidades de probabilidad.
Otro ejemplo pudiera ser la glucosa, que se parte en bajo, normal, prediabético y diabético. Categorías para valores
de la variable continua de glucosa en sangre que se parten en 4 categorías interpretables.

\begin{definicion}[Funciones]
  Es una regla que asigna a cada elemento de un conjunto inicial (Dominio) un único elemento de un conjunto 
  final (codominio). La denotamos como $f(x)=y$ donde $f:X\rightarrow Y$ denota cuál conjunto es el dominio ($X$) y 
  cuál el codominio ($Y$), de dónde a donde va y cómo asigna un elemento $x$ a uno $y$.
\end{definicion}

\section{Probabilidad}
Los conceptos previos más escenciales y necesarios de la teoría de probabilidad para tener nociones de 
lo que se realiza en estadística.\\
La escencia de la probabilidad es que si algo puede describirse en terminos probabílisticos $\Rightarrow$ podemos
analizarlo estadísticamente. La estadística por si sola es el mero análisis descriptivo de los datos más laxo posible,
para poder inferir requerimos de la probabilidad y para ello se requiere que la muestra obtenida se haya obtenida
a partir del azar. Sin Aleatoriedad en el muestreo, no hay nada de lo que se revisa a continuación.

\begin{definicion}[Axiomas de Kolmogrov]\label{def:Kolmogrov}
  Dado esp. muestral $\Omega$ y $\sigma$-álgebra (partición) de $\Omega$,
    una función de probabilidad es una función $P$ con dominio $\sigma$
        \begin{enumerate}
        \item $P(A)\geq 0$ $\forall A \in \sigma$
        \item $P(\Omega)=1$
        \item Si $A_{1},A_{2},... \in \sigma $ son disjuntos a pares, entonces 
        $P(\cup_{i=1}^{\infty}A_{i})= \sum_{i=1}^{\infty} P(A_{i})$
        \end{enumerate}
\end{definicion}

"frecuencia" es una forma particular de describir una probabilidad, la cual se entiende como la creencia en que un
evento pueda ocurrir. Definir la función de probabilidad es propio de un probabilista, no de un estadista. 
Es importante hacer el hincapié que teoría de probabilidad es distinto de la estadística. Esta última toma montones 
de resultados y herramientas de otros campos, existe en un terreno ambigüo entre la matemática pura y aplicada, mientras que
la teoría de probabilidad está más cerca de el área pura, en particular cuando se abarca desde teoría de la medida.

\begin{ejemplo}
  Ejemplo de la moneda \\
  $P(\conj{H})=P(\conj{T})$, mejor visto $P(\conj{H})+P(\conj{T})=1$ 
\end{ejemplo}

El tercer punto de \cref{def:Kolmogrov} muchas veces se cambia por: \\
$A \in \beta$ y $B \in \beta $ disjuntos, entonces $P(A \cup B)=P(A)+P(B)$

\begin{teorema}
  Si $P$ es probabilidad funcionan y $A \in \sigma \Rightarrow$
  \begin{enumerate}
    \item $P(\emptyset)=0$
    \item $P(A)\leq 1$
    \item $P(A^{c})=1-P(A)$
  \end{enumerate}
\end{teorema}

\begin{teorema}
  Si $P$ es una función de probabilidad y $A$ y $B$ son cualesquiera eventos 
    en $\sigma$
    \begin{enumerate}
        \item $P(B \cap A^{c})=P(B)-P(A \cap B)$
        \item $P(A\cup B) = P(A) + P(B) -P(A\cap B)$
        \item Si $A \subset B$ $\Rightarrow P(A) \leq P(B)$ 
    \end{enumerate}
\end{teorema}

\begin{proposicion}[Desigualdad Bayesiana]
  \[ P(A \cap B) \geq P(A)+P(B)\]
\end{proposicion}

\begin{teorema}
  Si $P$ es una función de probabilidad
        \begin{itemize}
            \item $P(A)=\sum_{i=1}^{\infty} P(A\cap C_{i})$ 
            $ \forall C_{1},C_{2},...$ partición de $\Omega$
            \item $P(\cup_{i=1}^{\infty}A_{i})\leq\sum_{i=1}^{\infty} P(A_{i})$ $\forall A_1, A_2,...$
        \end{itemize}
\end{teorema}

Los temas de conteo en la sección de probabilidad se decartaron en cierta medida durante el curso para
optimizar el tiempo en otros temas, aún así, se hicieron menciones breves particularmente para notar las diferencias 
en los métodos de conteo.\\
La manera más típica desde la cual entender o tener una noción de los métodos de conteo es con ejemplo de boleto
de lotería. Supugamos tenemos un boleto de lotería de seis números digamos que cada número es del 1 al 9,
entonces habría $9 \times 8 \times 7 \times 6 \times 5 \times 4$ posibles boletos. Si en su caso, los números 
fueran del 1 al 44, entonces habría $44 \times 43 \times 42 \times 41 \times 40 \times 39$.

\begin{table}[H]
  \centering
  \caption{Métodos de conteo}
  \begin{tabular}{@{}l >{\centering\arraybackslash}p{2.2cm}
                    >{\centering\arraybackslash}p{2.2cm}
                    >{\centering\arraybackslash}p{2.2cm}@{}}
    \toprule
    \textbf{Propiedad} & sin reemplazo & con reemplazo \\
    \midrule
    Sin orden           & $\frac{n!}{(n-r)!}$ & $n^{r}$ \\
    Con orden         & $\binom{n}{r}$ & $\binom{n+r-1}{r}$ \\
    \bottomrule
  \end{tabular}
\end{table}

\begin{teorema}
  Si un trabajo consiste en $k$ tareas que se puedan hacer de $n_j$ maneras, entonces
  el trabajo puede hacerse de $n_1 \times n_2 \times ... \times n_k$ formas.
\end{teorema}

\begin{observacion}
  Es importante tener en cuenta si están evaluando conteos en un proceso con o sin orden. Si tuvieramos
  que obtener las combinaciones de 4 que podemos formar con $\conj{2,4,9,12}$ sin que importe el orden tendríamos
  $\frac{12}{256}$ y con orden $\frac{1}{35}$. 
\end{observacion}

\section{Probabilidad condicional e independencia}
A veces hay que ampliar $\Omega$ con nueva información, para lo que hay que analizar el cálculo de probabilidad
o usar probabilidad condicional. 

\begin{ejemplo}
  Probabilidad de tener 4 aces (de maso de 52), puede entenderse desde la combinatoria
  $\binom{52}{4}$, pero puede también entenderse desde un enfoque de probabilidad condicional, es decir,
  pensando, que primero tienes 4 de 52, después tendrías 3 de 51; $\frac{4}{52} \times \frac{3}{51} \times \frac{2}{50} \times \frac{1}{49}$
\end{ejemplo}

\begin{definicion}
  Si $A$ y $B$ son eventos en $\Omega$ y $P(B)>0$ $\Rightarrow P(A|B)$\\
  Donde la probabilidad condicional de $A$ dado $B$ es $P(A|B)=\frac{P(A\cap B)}{P(B)}$\\
\end{definicion}

Notése que si $A\cap B = \emptyset \Rightarrow P(A\cap B)=0 \Rightarrow P(A|B)=0$

\begin{teorema}[Teorema de Bayes Generalizado]\label{thm:Bayes}
  Sea $A_{1},A_{2},...$ una partición de $\Omega$ y $P(B)>0$ con $B \subset \Omega$,
        entonces
            \[ P(A_{i}|B)=\frac{P(A_{i})P(B|A_{i})}{\sum_{j=1}^{\infty} P(B|A_{j})P(A_{j})}=\frac{P(A)P(B|A)}{P(B)}=\frac{P(A)P(B|A)}{P(B)P(B|A)+P(A^c)P(B|A^c)}\]
\end{teorema}

\begin{proof}
  Teorema de Bayes caso base\\
  $P(A|B)P(B)=\frac{P(A\cap B)}{P(B)}$\\
  $P(A|B)P(B)=P(A\cap B)$ de manera análoga para $P(B|A)$ tendríamos \\
  $P(B|A)P(A)=P(A\cap B)$ por lo que podemos igualar ambas expresiones\\
  $P(A|B)P(B)=P(B|A)P(A)$ que despejando, se obtiene
  \[ P(A|B)=P(B|A)\frac{P(A)}{P(B)}\]
\end{proof}

\begin{definicion}
  Dos eventos $A$ y $B$ son estadísticamente independientes si $P(A\cap B)=P(A)P(B)$.
  Si $B$ no tiene efecto en $A$ $\Rightarrow P(A|B)=P(A)$
\end{definicion}
Notése que en caso de independencia tendríamos $P(A)= \frac{P(B|A)}{P(B)}P(A)$ que es
equivalente a $1=\frac{P(B|A)}{P(B)}$ $\Rightarrow P(B)=P(B|A)$, lo cual es efectivamente 
$P(A \cap B)=P(B)P(A)$

\begin{definicion}
  Si $A$ y $B$ son independientes 
  \begin{enumerate}
    \item $A$ y $B^c$
    \item $A^{c}$ y $B$
    \item $A^{c}$ y $B^{c}$
  \end{enumerate}
  Son independientes.
\end{definicion}

\begin{proof}
  $P(A\cap B^{c})=P(A)-P(A\cap B)$\\
  $=P(A)-P(A)P(B)$\\
  $=P(A)(1-P(B))$\\
  $=P(A)P(B^c)$
\end{proof}

\begin{definicion}[Mutuamente independiente]
  Se dice que un evento es mutuamente independeinte cuando se cumple que 
  \[ P(\bigcap^k_{j=1}A_{ij})=\prod^{k}_{j=1}A_{ij} \]
\end{definicion}

\section{Funciones de Distribución}

\begin{definicion}[Variable Aleatoria]
  Una variable aleatoria es una función $f:\Omega \rightarrow \R$\\
  A esta función usualmente se le denota como $X$ cuando se habla de la función y $x_i$ 
  cuando se habla de la i-ésima observación. La minúscula para cuando es un valor observado, la mayúscula 
  es para la función.
  \[ X=x_i \iff \sigma_j \in \sigma | X(\sigma_j)=x_i\] donde $sigma$ es una partición de $\Omega$ y $\sigma_j$ un elemento de la partición.\\
  Una manera más clara de definirlo es expresándola expresamente como una probabilidad de eventos de las particiones 
  del espacio muestral.
  \[ P_X(X=x_i)=P(\conj{\sigma_j \in \sigma: X(\sigma_j)=x_i}) \]
\end{definicion}

\begin{ejemplo}
  Consideremos $X$ variable aleatoria como la cantidad de águilas que se pueden obtener en 3 volados, y sea 
  $\omega$ nuestro espacio muestral, en este caso sería todos los posibles resultados de hacer tres volados, donde
  $\sigma$ sería un único de estos resultados (nuestra partición del espacio muestral).
  Abreviemos águila como $A$ y sello como $S$.

  \begin{table}[H]
  \centering
  \label{tab:volados}
  \begin{tabular}{@{}lcccccccc@{}}
    \toprule
    $\sigma$ & AAA & AAS & ASA & SAA & ASS & SAS & SSA & SSS \\
    \midrule
    $X(\sigma)$ & 3 & 2 & 2 & 2 & 1 & 1 & 1 & 0 \\
    \bottomrule
  \end{tabular}
\end{table}

\end{ejemplo}


\begin{definicion}
  La función de distribucipon acumulada (cdf en inglés) de una variable aleatoria $X$ denotada por 
  $F_{X}(x)$ es 
  \[ F_X(x)=P_X(X\leq x) \forall x\]
\end{definicion}

\begin{ejemplo}
  Retomando el ejemplo de los 3 volados, donde $X=$águila puede describirse como

  $F_X(x)=\left\{
  \begin{array}{l}
    0 \quad -\infty \leq x > 0 \\
    \frac{1}{8} \quad 0 \leq x < 1 \\
    \frac{1}{2} \quad 1 \leq x < 2 \\
    \frac{7}{8} \quad 2 \leq x < 3 \\
    1 \quad 3\leq x < \infty
  \end{array}
\right.$
\end{ejemplo}

Pudimos organizar \cref{tab:volados} como 

\begin{table}[H]
  \centering
  \begin{tabular}{@{}lcccccccc@{}}
    \toprule
    $x$ & 0 & 1 & 2 & 3 \\
    \midrule
    $P_X(X=x)$ & $\frac{1}{8}$ & $\frac{3}{8}$ & $\frac{3}{8}$ & $\frac{1}{8}$\\
    \bottomrule
  \end{tabular}
\end{table}

Lo cual es un formato más intuitivo en el cual entender cómo se estructuró la función de probabilidad acumulada.
Y también la función propia de distribución de probabilidad que se revisará un poco más adelante
\begin{teorema}
  $F(x)$ es una función de probabilidad acumulada $\iff$
  \begin{itemize}
    \item $\lim_{x \to - \infty }F_{X}(x)=0$ y $\lim_{x \to + \infty }F_{X}(x)=1$
    \item $F_{X}(x)$ es no decreciente
    \item $F_{X}(x)$ es continua por la derecha
  \end{itemize}
\end{teorema}

Una variable aleatoria $X$ es continua $F_X(x)$ es función continua. $X$ es una variable aleatoria discreta si 
$F_X(x)$ es una función discreta o a pares.

\begin{ejemplo}[Caso Binomial]
  $P(X=x)=(1-p)^{x-1}p$\\
  $P(X\leq x)=\sum_{i=1}^{x}P(X=i)=\sum_{i=1}^{x}(1-p)^{i-1}p$
  $\Rightarrow F_X(x)=1-(1-p)^{x}$
\end{ejemplo}
Este caso concreto puede corrobarse con el ejercicio de la moneda y volados y lo que se había presentado
anteriormente. Recordando claro que se uso una moneda balanceada, es decir, con $p=\frac{1}{2}$.

\begin{definicion}
  Una variable aleatoria $X$ y $Y$ está idénticamente distribuidas si $\forall A \in \sigma$
  $P(X\in A)=P(Y\in A)$.
\end{definicion}

\begin{teorema}
  $X$ y $Y$ variables aleatorias se distribuyen idénticamente $\iff F_X(x)=F_Y(x) \quad \forall x$ 
\end{teorema}

Varias de las pruebas estadísticas y teoremas juegan directa o indirectamente con la probabilidad acumulada para 
analizar el tipo de distribución como pudiera serlo Kolmogrov-Smirnoff o $\chi^{2}$. Sin embargo, este tipo de 
probabilidad puede ser insuficiente para describir otros aspectos de mucho mayor interés o describirlos puede ser 
mucho más complicado, para eso se constuye otra serie de conceptos complementarios.

\begin{definicion}
  La función de probabilidad de masa es discreta y está dada por 
  \[ f_X(x)=P(X=x) \]
\end{definicion}
Está probabilidad se refiere al 'tamaño del salto' de $F_X(x)$ en este caso.\\
El ejemplo más clásico es el de la probabilidad geométrica\\
$f_X(x)=P(X=x)=\left\{
  \begin{array}{l}
    (1-p)^{x-1}p \quad x\in \Z^{+}\\
    0
  \end{array}
\right.$\\

$f_X(x)$ es entonces el tamaño de de salto de la distribución de probabilidad acumulada
de la variable aleatoria $X$.\\
$P(a\leq x \leq b)=\sum_{k=a}^{b} f_X(k)=\sum_{k=a}^{b}(1-p)^{x-1}p$\\
$P(x \leq b)=\sum_{k=1}^{b} f_X(k)=F_X(b)$\\

Es fácil ver cómo en este caso funciona como un salto y al observar el cómo sumar los valores hasta un punto 
nos da la acumulada hasta ese punto nos invita a pensar en uno de los resultados más claves. Pareciese que la
distribución de masa y acumulada en un caso continuo pudieran acoplarse al teorema fundamental del cálculo,
que revisando los límites la clave está en revisar.\\

$P(X=x)\leq P(x-\epsilon < X \leq x)=F_x(x)-F_X(x-\epsilon)$\\
$P(X\leq x)=F_X(x)=\int_{-\infty}^{\infty}f_X(t)dt$\\
$\frac{d}{dx}F_X(x)=F_X(x)$ por T.F.C.

\begin{definicion}[\termino{Función de Densidad de Probabilidad}]
  $f_X(x)$ de una variable aleatoria continua $X$ es la función que satisface
  \[ F_X(x)=\int_{-\infty}^{\infty}f_X(t) dt \quad \forall x\]
\end{definicion}
De manera análoga en el caso discreto debe cumplirse para la sumatoria en lugar de la integral.\\
Además de esto, de los axiomas de probabilidad e inferencia detrás de la construcción de la función de 
densidad, podemos intuitivamente entender.
\begin{teorema}
  $f_X(x)\geq 0 \quad \forall x \quad \iff$
  \begin{itemize}
    \item $f_X(x) \geq 0 \quad \forall x$
    \item $\sum_{x}f_X(x)=1$
    \item $\int_{-\infty}^{\infty}f_X(x)dx=1$
  \end{itemize}
\end{teorema}
Estos resultados son análogos en el caso discreto, prácticamente solo cambiando integral por sumatoria.\\
Uno de los conceptos más importantes son las transformaciones, prácticamente todo se fundamente ahí. En general,
uno de los conceptos más importantes en matemática es el de la 'invarianza', qué es revisar cuánto podemos deformar 
o transformar algo conservando sus propiedades o cuáles son las que permanecen tras una cierta transformación.\\
El por qué queremos transformar yace en las facilidades o practicidades que esto brinda. La más conocida es la
normalización o estandarización de las variables que transforma la distribución normal en una cierta escala a una
normal estándar de media 0 y desviación 1. Pero cómo está hay montones de transformaciones que sean implícita o 
excplícitamente.\\
Si $X$ es una variable aleatoria con $F_X(x)\Rightarrow$ toda función $g(x)=y$ de $X$, es también una variable aleatoria 
que podemos describir en términos de $X$ para todo evento $A$ del espacio muestral $\Omega$. Donde normalmente condisderamos
$g(x):x\rightarrow y $ como $P(g(x)\in A)=P(y\in A)$ donde $g^{-1}(\conj{y})=\conj{x\in X: g(x=y)}$.\\
$P(g(x)\in A)=P(X\in g^{-1}(A))$\\
$f(g(y))=P(g=y)=\sum_{x\in g^{-1}(y)}P(X=x)$\\
$\quad =\sum_{x\in g^{-1}(y)}f_X(x) \quad y\in Y$.\\

Una de las transformaciones con las que indirectamente se juega mucho y que sirve mucho para teoría y otras 
cuestiones es la uniforme, que es aquella donde todo tiene la misma probabilidad de ocurrir. 

\begin{figure}[H]
  \centering
  \begin{tikzpicture}[scale=1.2]
    \draw[->] (-0.3,0) -- (3,0) node[right] {$x$};
    \draw[->] (0,-0.3) -- (0,2.5) node[above] {$y$};
    \draw[domain=0:3, smooth, thick, blue] plot (\x, 1);
  \end{tikzpicture}
  \caption{Gráfico de función de masa uniforme}
\end{figure}

Un ejercicio útil es el de dibujar la función de probabilidad acumulada de esta función.

\begin{ejemplo}[tranformación logarítimica]
  Sea $g(X) = -log(X)$\\
  $\frac{d}{dx}g(x)=\frac{d}{dx}(-log(x))=\frac{-1}{x}<0 \quad \forall x \in (0,1)$\\
  $g$ entonces es decreciente. Además notemos que el rango de $X$ es $[0,1]$ y el de 
  $-log(x)$ es $(0,\infty)$\\
  para $y=-log(x), \quad y>0$ $x=e^{-y}\Rightarrow g^-1(y)=e^{-y}$\\
  $F_Y(y)=1-F_X(g^-1(y))=1-F_X(e^{-y})=1-e^{-y}=F_X(y)=x$
\end{ejemplo}
Es un ejemplo de lo que ocurre con las funciones probabilidad acumulada en el caso concreto de 
la transformación logarítimica que es una de las que al final terminará siendo de las más usados por 
temas de escalabilidad y ajuste de distintas dimensiones. Este ejemplo también sirve para ilustrar 
más útilidades de un análisis de derivadas al momento de evaluar o considerar tranformaciones. 
Para el uso típico de pruebas no se hace porque ya se tienen varios resultados, pero si se busca entender
el fondo y apreciar resultados así como revisar los teoremas de construcción de las principales distribuciones
y tipos de variables.

\begin{teorema}
  Inducir $f_Y(y)$ dado $f_X(x)$ es continua y $g^{-1}(y)$ tiene derivada continua en $Y$. 
  \[f_y=\left\{
  \begin{array}{l}
    F_X(g^{-1}(y)) | \frac{d}{dy} g^{-1}(y) | y\in Y \\
    0 \quad e.c.o.c.
  \end{array}
\right.\]
\end{teorema}

\begin{ejemplo}[Tranformación Cuadrada]
  Suponga que $X$ es una variable aleatoria continua para $y>0 \Rightarrow$ la función de 
  acumulación $Y=X^2$ es\\
  $F_Y(y)=P(y\leq Y)=P(X^2\leq Y)=P(-\sqrt{y}\leq X \leq \sqrt{y})=F_X(\sqrt{y})-F_X(-\sqrt{y})$ sería entonces la función 
  de probabilidad acumulada.\\
  $f_Y(y)=\frac{1}{2\sqrt{y}}(f_X(\sqrt{y})+f_X(-\sqrt{y}))$ integrando lo anterior.
\end{ejemplo}
Este ejemplo particular es de importancia porque la distribución ji-cuadrada viene precisamente de elevar una normal estándar 
al cuadrado, que es el ejemplo a continuación. Ese hecho es de vital importancia para entender la mayoría de fundamentos de las pruebas 
paramétricas, ya que se construyen realizando operaciones sobre normales, las cuales van consturyendo las dsitribuciones especiales
como la ji-cuadrada, t-student, F-Fisher.

\begin{ejemplo}[Relación normal-$\chi^{2}$]
  Sea $X$ la distribución normal, entonces su densidad está dada por 
  $f_X(x)=\frac{1}{\sqrt{2\pi}}e^{-\frac{x^2}{2}}$ con $-\infty<x<\infty$\\
  Consideremos $g(x)=x^2$, sabemos que es monótona en $(-\infty,0),(0,\infty)$, así que consideremos $Y=(0,\infty)$\\
  $\Rightarrow f_Y(y)=\frac{1}{\sqrt{2\pi}}e^{\frac{-(-\sqrt{y})^{2}}{2}}(|\frac{-1}{1+2\sqrt{y}}|+|\frac{1}{2\sqrt{y}}|)$\\
  $\quad = \frac{1}{\sqrt{2\pi}} \frac{1}{\sqrt{y}}e^{-\frac{y}{2}}$ en $0<y<\infty$\\
  Esa función de densidad de probabilidad es la misma que la de una ji-cuadrada de un grado de libertad.
\end{ejemplo}

Otro resultado escencial para entender las distintas pruebas estadísticas es sobre la distribución uniforme,
que ya se había mencionado que funciona como una distribución de 'transición' durante algunos procesos.

\begin{teorema}
  Sea $X$ variable aleatoria con función de probabilidad acumulada $F_X(x)$ y define $Y$ como
  $Y=F_X(x)$ $\Rightarrow Y$ se distribuye uniformemente en $(0,1)$. Es decir, 
  $P(Y\leq y)=y$ para $0<y<1$.
\end{teorema}

\chapter{Estimación}

\section{Valor Esperado}

\begin{definicion}
  \[E(g(x))=\left\{
  \begin{array}{l}
    \int_{- \infty}^{\infty}g(x)f_{X}(x)dx \quad X \cong \R\\
    \sum_{x\in X}g(x)f_{X}(x) \quad X \cong \N
  \end{array}
\right.\]
\end{definicion}

El valor esperado es tal vez el concepto más importante de la teoría de probabilidad que 
la estadística inferencial hace uso. Se puede entender a la estimación como ese un análisis de valores 
esperados. La estimación puntual no es tan usada en investigación porque no es algo que se preste a una 
interpretabilidad clara, pero su formato como estimación de intervalo que posee niveles de confianza y se 
alinea apropiadamente con las pruebas de hipótesis. 

\begin{observacion}
  si $E(g(x)) = \infty \Rightarrow \not \exists E(g(x))$ 
\end{observacion}

Conocer las propiedades del valor esperado es de vital importancia para entender cualquier texto en estadística 
o en cualquiera de sus ramas. Este es a fin de cuenta el concepto del que se establece la media, varianza, covarianza y 
todos los conceptos que se discuten en cualquiera de las ramas y nichos específicos de la estadística, además de que sus
modificaciones son a través de el esperado, modificaciones entendidos como la función media o función de covarianza en 
análisis de series de tiempo, por mencionar alguno. 

\begin{teorema}
  Sea $X$ una variable aleatoria y sean $a,b$ y $c$ constantes $\Rightarrow g_1(x)$ y $g_2(x)$ con 
  $E\in \R$.
  \begin{itemize}
    \item $E(ag_1(x)+bg_2(x)+c)=aE(g_1(x)+bE(g_2(x))+c)$
    \item Si $g_1(x)\geq 0 \quad \forall x \Rightarrow E(g_1(X))\geq 0$
    \item Si $g_1(x)\geq g_2(x) \quad \forall x \Rightarrow E(g_1(x))\geq E(g_2(x))$
    \item Si $a \leq g_1(x)\leq b \quad \forall x \Rightarrow a \leq E(g_1(x))\leq b$
  \end{itemize}
\end{teorema}

\section{Estimador}
Previamente se discutió que el fundamento de la estadística inferencial es la estimación, el buscar describir una población 
con una muestra. Estas estimaciones pueden ser puntuales como la media aritmética, o pueden ser en intervalos que cuando 
trabajamos con hipótesis u otros factores, se convierten en intervalos de confianza. Cada tipo de estimación tiene su función 
y sus matices. \\
Un estimador puntual es una función $W(X_1,...,X_n)$ de una muestra. Cualquier estadístico en un estimador puntual.
Usualmente se denota $W(x_1,...,x_n)$ como un valor estimado, mientras que $W(X_1,...,X_n)$ es un estimador, una función.
Muchos estimadores se obtuvieron de manera intuitiva primero y después se vió que encajaban con la teoría propuesta, de hecho, 
hay quienes recomiendan la intuición como la manera más eficiente de conseguir un estimador, claro que el problema se vuelve 
el fundamento que le respalde, pero tal vez revisar que cumpla con condiciones es más sencillo que construirlo a partir de ellas.\\
Como aquí ya comenzamos a hablar de manera mucho más formal sobre la muestra vale la pena remarcar el concepto.
Las variables aleatorias $X_1,X_2,...,X_n$ son llamadas muestras aleatorias de tamaño $n$ de la población
$f(x)$ si $X_1,X_2,...,X_n$ son variables aleatorias mutuamente independientes y la fdp o fmp para cada $X_i$
es la propia $f(x)$. De manera alterna, $X_1,X_2...,X_n$ son llamadas variables aleatorias independientes e idénticcamente 
distribuidas por $f(x)$, comúnmente abreviado como iid.

\section{Método de Momentos}

Los estimadores más usuales, comúnes y básicos vienen del método de momentos, además de que este método es pilar para todas 
las disciplinas adyacentes de la estadística. Todo que tome algo de la estadística termina trabajando con el método de momentos
así que conocerle y entender sus generalidades puede ser muy fructífero. 

\begin{definicion}
  Para cada entero $n$ el $n$-ésimo momento de la variable aleatoria $X$ o $F_{X}(x)$ ($\mu_{n} '$), es
  \\ $\mu_{n} '=E(X)^n$ 
  Mientras que el $n$-ésimo momento central de $X$, $\mu_{n} $, es
  \\ $\mu_{n}=E(X-\mu)$,
  \\ donde $\mu = \mu_{1} '=E(X)$
\end{definicion}

El método de momentos a nivel práctico se realiza planteando y resolviendo el siguiente sistema de ecuaciones:
\[m_1=\frac{1}{n} \sum_{i=1}^{n} X_{i}^{1} \quad \mu_1' =E(X^1) \]
\[m_2=\frac{1}{n} \sum_{i=1}^{n} X_{i}^{2} \quad \mu_2' =E(X^2) \]
\[...\]
\[m_k=\frac{1}{n} \sum_{i=1}^{n} X_{i}^{k} \quad \mu_k' =E(X^k) \]

$\mu_j'$ es una función de $(\theta_1,...,\theta_k)$ el estimador $(\hat{\theta_1},...,\hat{\theta_k})$
de $(\theta_1,..,\theta_k)$ se obtiene resolviendo para $(\theta_1,...,\theta_k)$ en términos de $(m_1,...,m_k)$
\[ m_1=\mu_1'(\theta_1,...,\theta_k) \]
\[ m_2=\mu_2'(\theta_1,...,\theta_k) \]
\[...\] 
\[ m_k=\mu_k'(\theta_1,...,\theta_k) \]

La media es el primer momento de centralidad, entendido también como el centro de masa de la distribcuión de probabilidad.
La varianza de una variable aleatoria X es el segundo momento central, $VarX=E(X-EX)^2$. 
La raíz cuadrada positiva de $VarX$ es la desviacion estándar de $X$ cuando hablamos de una dsitrubción normal.
Cuando la media tiende a valores grandes, $X$ es más variable, pero sí $VarX=E(X-EX)^2=0$,
entonces $X=E(X)$ con probabilidad 1 y por tanto la varianza es 0.

\begin{ejemplo}
  Revisemos cómo se utiliza el método de momentos para hallar los estimadores de los primeros dos momentos (media y varianza) 
  en una distribución normal. Así que tomemos $X_1,...,X_n~N(\mu,\sigma^{2})$\\
  para ello plantemos las ecuaciones 
  \[m_1=\frac{1}{n}\sum X_i \quad \mu_1'=\mu\]
  \[m_2=\frac{1}{n}\sum X_i^2 \quad \mu_2'=\mu^2+\sigma^2\]
  Tendríamos entonces $\frac{1}{n}\sum X_i= \mu$\\
  $\frac{1}{n}\sum X_i^2-(\frac{1}{n}\sum X_i)^2$\\
  $\quad = \frac{1}{n}\sum X_i^2-\frac{1}{n^2}\sum X_i^2=\frac{1}{n}(\sum X_i^2-\frac{1}{n}\sum X_i^2)$\\
  $\quad = \frac{1}{n}\sum X_i^2-\bar{X}^2=\sigma^2$
\end{ejemplo}
De aquí se puede apreciar que la fórmula para estimar varianza poblacional y media vienen o pueden estar dadas por el 
método de momentos, el cual como se ha repetido durante varios conceptos, es fundamental para entender cómo funciona la 
estadística, más aún, cómo funcionan las ramas adyacentes de esta. De aquí también se pueden intuir propiedades de la varianza
como que si $X$ es una variable aleatoria con varianza finita, entonces para cualesquiera constantes $a$ y $b$; $Var(aX+b)=a^{2}Var(X)$.\\

Resolver el sistema anterior puede ser demasiado trabajo, y tiene un cierto patrón que se desglozó para desarrollar 
un concepto muy particular que es la función generatriz de momentos.\\

\begin{definicion}
  Sea $X$ una variable aleatoria con dpa $F_{X}$. La función genartriz de momentos (fgm) de 
  $X$ o $F_{X}$, denotada por $M_{X}(t)$, es
  \[M_{X}(t)=E(e^{tX})\] 
\end{definicion}


\section{Máxima Verosímilitud}



El rango de MLE es el mismo del parámetro. Pero cabe resoltar que este método posee 3 problemas muy sensibles:
\begin{itemize}
  \item Maximización/ Convexidad
  \item Sensibilidad Númerica
  \item Solo acepta soluciones únicas
\end{itemize}

Las específicaciones, trabajo y ejemplos clave de este método fueron omitidos. Muchos detalles importantes 
algo ignorados ya que hay que tener nociones algo concisas de logartimos, cotas, piatorias, cálculo vectorial, 
además de una fluidez en el maneho del supremo para saber apreciar sus propiedades.\\

Una de sus mayores virtudes es la invarianza ante tranformaciones. Por ejemplo si $\hat{\theta}$ es MLE de $\theta$,
entonces $f(\hat{\theta})$ es MLE de $f(\theta)$. Esto también afecta a su varianza así que tiene sus detalles, por ejemplo,
considermos $\theta = \mu = \frac{1}{n}\sum X_i\Rightarrow \quad \theta^2 = \mu^2 = (\frac{1}{n}\sum X_i)^2$.

MLE $\hat{\theta}$normalmente se obtiene a través de 
\[ \sum_{i=1}^{n} \frac{\delta ln(L(x_i|\theta))}{\delta \theta}=0 \]
\[ \frac{\delta^2 ln(L(x_i|\theta))}{\delta \theta}{\delta \theta^2} \hat{\theta}<0 \]
El ejemplo más clásico 

\begin{ejemplo}
  Recibes 4 cartas, pero 1 se extravía. Si de las 3 que sí tienes, 2 tienen facutras/recibos y la otra dinero/premio
  ¿cuál sería un buen estimado del número de facturas de las 4 cartas? \\
  En este caso hay dos opciones solamente, o 2 son facturas (las que tenemos), o la que perdimos también lo es y deberíamos 
  tener 3 facturas. Para esto no pensamos en lo probable que fuese recibir una factura o un premio, intentamos pensar desde
  el enfoque cuál es la probabilidad de cada cosa para perderse. En este caso asumamos que todas tienen la misma probabilidad de 
  perderse, así tendríamos que\\
  $\frac{\binom{2}{2} \binom{2}{1}}{\binom{4}{3}}=\frac{1}{2}$ para cuando asumimos que son solo 2 facutras.
  $\frac{\binom{3}{2} \binom{1}{1}}{\binom{4}{3}}=\frac{3}{4}$ para cuando asumimos que son solo 3 facutras.
\end{ejemplo}
  La probabilidad que nos da mayor es si pensamos que hay 3 facturas, esto es porque ya tenemos dos de nuestras tres cartas 
  como facturas. Por eso, el estimador de máxima verosímilitud sería que hay una probabilidad de $\frac{3}{4}$ de que la otra carta
  sea una factura. El punto aquí es ver que hay dos posibles estimadores, pero máxima verosímilitud es uno concreto, a partir de un 
  enfoque concreto. 


\section{Bayesiano}
En nuestra estadística normalmente el parámetro $\theta$, es desconocido, pero es un valor fijo, es cantidad. Observamos variables aleatorias
y obtenemos información del parámetro $\theta$. El enfoque Bayesiano varía $\theta$ según su distribución (priori) la cuál es subjetiva.
Se toma un a muestra y actualizamos la distribcuión a posteriori. \\
Usalmente se denota la priori como $\pi(\theta)$, mientras que muestro es 
$f(x|\theta)$. Donde se denota $m(x)$ como la probabilidad marginal\\
$m(x)=\int f(x|\theta)\pi(\theta)d\theta$\\
Donde la distribución es $\pi(\theta | x)=\frac{f(x|\theta)\pi(\theta)}{m(x)}$\\
Notése que $\pi(\theta|x)$ depende de haber observado la muestra.

\begin{ejemplo}
  ---FCFM vs FACPYA---
\end{ejemplo}

\section{Propiedades de Estimador}

\begin{ejemplo}
  En una distribución binomial, la media está estimada por $n\theta$, es decir, $E(X)=n\theta$.\\
  Busquemos probar que $\frac{x}{n}$ es insesgado.\\
  $E(\frac{x}{n})=\frac{1}{n}E(x)$
  $\quad =\frac{1}{n}n\theta$
  $\quad = \theta$, por lo tanto, $\frac{x}{n}$ es insesgado. 
\end{ejemplo}

Sobre consistencia existe un resultado que dice 

\begin{teorema}
  Si $\hat{\theta}$ es estimador insesgado de $theta$ y $var(\hat{\theta})\rightarrow 0$ cuando $n\rightarrow \infty$, 
  entonces $\hat{\theta}$ es estimador consistente de $\theta$.
\end{teorema}

Se entiendo como que al aumentar la muestra, su varianza se reduzca siendo ya un estimador insesgado. Un estimador que no
fuera consistente podría no reducir su varianza a pesar de tener más muestra o incluso empeorar su varianza. Suena a contraproducente
Pero hay estimadores que no cumplen con esto a pesar de que es algo con lo que siempre se busca tener.\\
Uno de los teoremas más importantes en esta sección que no se habló por ser algo más técnico es el 
Fisher-Neyman. Resultado importante con rescpecto a mínima varianza, ya que en la discusión al comparar 
$\frac{Var(\hat{\theta_1})}{Var(\hat{\theta_2})}$ (que es una $\chi^2$).

\begin{definicion}
  Los cuadrados medios del error (CME) de un estimador $W$ de un parámetro $\theta$, es la función definida como:
  \[E_\theta(W-\theta)^2 \] 
\end{definicion}
Los CME es de los conceptos más importantes de estadística, ya que a partir de aquí se construyen bastantes de los diseños más 
utilizados en el análisis de varianza e incluso en regresiones. Conceptos semejantes son análogos en campos adyacentes a la 
estadística. 
\begin{observacion}
  Es importante notar que cualquier función creciente de $|W-\theta|$ es una buena medida del $performance$ de la medició,
  además\\
  $E_\theta(W-\theta)^{2}=Var_\theta(W)+(E_\theta(W-\theta)^2)^{2}$\\
  $\quad = Var_\theta(W) + (Bias_\theta (W))^{2}$
\end{observacion}

Pensemos en $X_1,X_2,...,X_n$ variables aleatorias iid $N(\mu, \sigma^{2})$, donde sabemos que $E(\bar{X})=\mu$ y $E(S^{2})=\sigma^2$ son
estimadores de media y varianza respectivamente, así que midiendo los CME obtendríamos que\\
$E(\bar{X}-\mu)^{2}=Var(\bar{X})=\frac{\sigma^{2}}{n}$ que se puede interpretar como que el error en la estimación de la media depende 
directamente de la varianza de la población y es inversamente proporcional al tamaño de la muesta, a mayor muestra, menor será CME de la
media estimada.\\
$E(\bar{X}-\sigma^{2})^{2}=Var(S^2)=\frac{2\sigma^{4}}{n-1}$ el cual tiene una intepretación muy similar al anterior, solo que el error
crece bastante más por unidad de varianza, así que se requiere muchísima más muestra. De aquí sale la popular expresión de 
$\frac{n-1}{n} S^{2}$.\\
Es importante notar que los estimadores no se 'controlan' bajando el sesgo. Si el sesgo bajo, la varianza va a aumentar, lo que 
aumenta los CME. Esto sirve para ilustrar que las maneras de evaluar un estimador abren el paronoma a un problema vasto como seleccionar
un estimador, para los cuales se construyen soluciones adecuadas. En el caso concreto de este ejemplo se creo el \textit{Coeficiente de Información}
como una cota para la mínima varianza del mejor estimador insesgado. Entre mayor sea el número de datos, más 'información' del parámetro 
tenemos y consecuentemente, una cota mucho más estrecha.\\
Esta sección generalmente continúa con un análisis detallado de la desigualdad Cramér-Râo, pero es un tema que se deja en reserva para 
un curso de estadística más formal. Solo es importante notar que es algo de importancia fundamental en la selección y construcción de
estimadores, que posee generalizaciones para funcionar sin importar de si se trata de un caso discreto o continuo. Además de que es una 
desigualdad que no necesariamente se satisface.

\chapter{ANOVA}

\section{Bases}

ANOVA estima medias de varias poblaciones a veces con hipótesis de normalidad. 
El concepto escencial de ANOVA yace en la pregunta ¿Cómo consigo la mayor cantidad de información en la 
mayor cantidad de poblaciones con la menor cantidad de observaciones?\\
Para resolver la pregunta se plantea un modelo similar al que se utiliza en regresiones; 
$y_{ij}=\theta_i + \epsilon_{ij}$ con $i=1,...,k$ y $j=1,...,n_i$ comunmente entendido como $k$ tratamientos 
y $n_i$ como las observaciones del $i$-ésimo tratamiento. 

\begin{table}[H]
  \centering
  \caption{Ordenamiento usual de observaciones en ANOVA}
  \begin{tabular}{@{}l >{\centering\arraybackslash}p{2.2cm}
                    >{\centering\arraybackslash}p{2.2cm}
                    >{\centering\arraybackslash}p{2.2cm}@{}}
    \toprule
    $T_1$ & $T_2$ & $...$ & $T_k$ \\
    \midrule
   -$y_{11}$ & $y_{21}$ &  $...$ & $y_{k1}$\\
    $y_{12}$ & $y_{12}$ &  $...$ & $y_{k2}$\\
    $...$ &  $...$ &  $...$ & $...$ \\
    $y_{1n_1}$ & $y_{2n_2}$ & $...$ & $y_{kn_k}$ \\
    \bottomrule
  \end{tabular}
\end{table}

\begin{observacion}
  En el modelo se asumen los errores como normales porque si $\epsilon_{ij}~N(0,1) \Rightarrow$\\
  $E(Y_{ij})=E(\theta_i + \epsilon_{ij})$\\
  $=E(\theta_i)+E(\epsilon_{ij})$\\
  $=E(\theta_i)$\\
  $\therefore E(y_ij)=E(\theta_i)$\\
  Donde $\theta_i$ son las medias de los tratamientos, que se presentan más claramento como $y_{ij}=\mu+\tau_i+\epsilon_{ij}$
\end{observacion}
  De $y_{ij}=\mu+\tau_i+\epsilon_{ij}$, $\mu$ es la media general de los efectos, $\tau_i$ es el efecto de los tratamientos. 
  La desviación es constante por los tratamientos. No podemos estimar $\mu$ y $\tau_i$ por seprado ya que se irrumpe por un problema de 
  'indeticafibilidad'. No sé si sea una traducción válida que exista en el español, pero se refiere a un problema dónde no se pueden 
  identificar, ergo, separar las estimaciones de manera adecuada. 

  \begin{definicion}[Identificable]
    Un parámetro $\theta$ de $\conj{f(x|\theta): \theta \in \Theta}$ es identificable si distintos valores de $\theta$ corresponder 
    a distintas funciones de densidad de probabilidad o funciones de masa de probabilidad. Es decir, si 
    $\theta \not = \theta '\Rightarrow f(x|\theta)\not = f(x|\theta ')$ para el mismo $x$.
  \end{definicion}
Los problemas de identificabilidad está determinada por el modelo, no los parámetros. Si no es identificable es complicado distinguir
y hacer inferencia porque si $f(x|\theta)=f(x|\theta ') \Rightarrow$ cualquier observación se ve igual, así que no sabemos si $\theta$
o $\theta '$ es el parámetro original de la misma población generan la misma función de máxima verosímilitud. 
Hasta este punto todo modelo era identificable, una normal hecha con una media y varianza que podemos distiguir. 
El modelo del ANOVA es muy intuitivo pero no es identificable.
\begin{observacion}
  Una manera de entender cuál es el origen del problema de identificabilidad revisando $y_{ij}=\mu + \tau_i + \epsilon_{ij}$ 
  hay $k+1$ parámetros pero solo hay $k$ medias, donde podemos entender la estimación de las $k$ medias de los $k$ tratamientos
  $E(Y_{ij})$ para $i=1,...,k$, entonces existe un elemento de $(\mu, \tau_1, ..., \tau_k)$ con distribucion repetida porque son 
  las medias de tratamientos y la población general dentro del modelo, por eso $\sum_{i=1}^{k} \tau_i=0$ es una hipótesis para garantizar
  la estimación de $k$ parámetros
\end{observacion}

\begin{observacion}[Hipótesis ANOVA]
  El modelo base del análisis está dado por
  \[ y_{ij}=\theta_i + \epsilon_{ij} \]

  \begin{itemize}
    \item $\epsilon_{ij}~ N(0, \sigma^2)$ donde $\sigma < \infty$
    \item $\epsilon_{ij}$ son independientes; $cov(\epsilon_{ij},\epsilon_{i'j'})=0 \quad \forall i, i', j, j'$ a menos que $i=i' \wedge j=j'$
    \item Homogeneidad de varianzas $\sigma_i^2=\sigma^2 \quad \forall i$
  \end{itemize}
\end{observacion}

Sin independencia y normalidad solo podríamos hacer estimación puntual que minimice varianza. siempre
que tengamos tamaños de muestra razonables y poblaciones no muy asimétricas, se puede llegar a hacer uso del 
CLT. Para poblaciones no normales pueden llegar a hacerse uso de transformaciones y cosas no tan sencillas para
poder seguir usando este método.\\
La homogeneidad de varianzas garantiza el poder hacer uso de transformaciones no lineales para evitar 
que no-violen las condiciones ANOVA. Por ejemplo, las transformaciones Box-Cox también garantizan robustez, 
aparte de lo que se viene hablando.

\begin{table}[H]
  \centering
  \caption{Hipótesis ANOVA}
  \begin{tabular}{@{}l >{\centering\arraybackslash}p{2.2cm}
                    >{\centering\arraybackslash}p{2.2cm}
                    >{\centering\arraybackslash}p{2.2cm}@{}}
    \toprule
    $H_0$ & $H_1$ \\
    \midrule
   -$\theta_i=0 \quad \forall i$ & $H_1=\exists i \not = j \quad | \theta_i \not = \theta_j$ \\
    $\theta_i=\theta_i \quad \forall i$ & \\
    $\theta_1=\theta_2=...=\theta_k$ & \\
    \bottomrule
  \end{tabular}
\end{table}

El interés de una prueba ANOVA es medir el efecto, muchas veces solo si existe.\\
Existe algo llamado 'contrastes' donde las hipótesis ANOVA pueden expresarse en 
términos de combinaciones lineal para entender mejor las comparaciones o elaborar 
hipótesis más específicas. Se usan empleando los conceptos clásicos de combinaciones lineales 
$\sum a_it_i \quad | \quad \sum a_i=0$. Va muy de la mano con el enfoque unión-Intersección
para diseño de hipótesis que tampoco se curbió. Hay un mundo de detalle en este campo pero no fue
vital, normalmente con entender el concepto de particionar un conjunto basta.\\



%% ============================================================
%%  ÍNDICE ANALÍTICO
%% ============================================================
\printindex

\end{document}
